\begin{abstract}
We study the population oracle-value curve indexed by treatment capacity for a binary treatment and outcome on a known discrete covariate alphabet. Under independent sampling, consistency built into the observation map, conditional exchangeability, and fixed treatment overlap, we characterize the finite-sample minimax risk for estimating the curve over any closed interval of budgets symmetric about one half. For every sample size \(n\geq1\), alphabet size \(d\geq2\), and interval endpoint \(b_0\in[0,1/2]\), the all-procedure expected squared supremum risk on \([b_0,1-b_0]\) is of order \(\min\{1,d/[n\log(ed)]\}\), with constants depending only on overlap. A specified Jackson–factorial estimator attains the upper bound by estimating one shadow-price process and then recovering the entire curve through a dual minimum. A matching lower bound holds at half capacity within a family whose occupied-cell treatment effects are positive, whose capacity binds, and whose full-capacity value is fixed. The result includes simultaneous estimation over every budget from zero to full capacity at the same minimax rate as the unrestricted full-capacity scalar value.
\end{abstract}

\section{Introduction}

A treatment capacity turns population value into a curve: each budget specifies the largest expected outcome attainable by a feasible assignment policy. Estimating that curve addresses questions about the value of additional capacity while allowing the most valuable recipients to change with the budget. This paper characterizes the finite-sample precision of that population target for binary treatment and outcome on a known discrete covariate alphabet that may grow with sample size.

Let \(d\) be the known number of covariate cells, \(n\) the number of independent observed records, and \(\epsilon\) a fixed treatment-overlap parameter. At budget \(b\), the oracle value \(V_b(P)\) is the maximum population outcome under policies whose average treatment share is at most \(b\), for a causal law \(P\) satisfying consistency, conditional exchangeability, and overlap. We measure estimation error by the expected squared largest error over \(I=[b_0,1-b_0]\), the budget interval with endpoint parameter \(b_0\in[0,1/2]\), and compare all measurable curve estimators. For every \(n\geq1\), \(d\geq2\), and \(b_0\in[0,1/2]\), \cref{obj:thm:matched-curve-frontier} establishes the risk order
\[
\min\left\{1,\frac{d}{n\log(ed)}\right\},
\]
with comparison constants depending only on \(\epsilon\). The statement includes the full interval \([0,1]\) and the single half-budget value.

The upper bound uses a dual representation of each budget value as the minimum of a common shadow-price process plus the capacity charge, as shown in \cref{obj:prop:dual-representation}. A pilot-local Jackson polynomial and factorial-moment statistics estimate that process across shadow prices; conditional averaging yields a curve estimator based on the observed sample in \cref{obj:def:jf-estimator}. The dual-representation proposition, \cref{obj:prop:dual-representation}, gives the factor-one contraction from uniform process error to curve error; \cref{obj:thm:curve-upper} gives the resulting estimator risk guarantee. This contraction lets estimation of the whole curve attain the minimax rate of the unrestricted full-capacity scalar value. The guarantee applies across the stated causal law class, including null cells, unequal occupied-cell masses, boundary outcome means, and treatment-effect ties.

The matching converse comes from a paired causal family. At half capacity, every occupied-cell treatment effect lies between \(1/8\) and \(3/8\), the available capacity is used, and the full-capacity value equals \(1/2\). In this family, the target obeys the affine identity \(V_{1/2}(P)=1/8+F_P(1/4)\), where \(F_P\) is the dual threshold process. An exact map sends paired samples from two unknown discrete distributions to observations from this causal family. Thus, an estimator of the half-budget value would yield an estimator of their \(L_1\) distance, giving the lower bound from \citet[Theorem 3 and the unknown-both-distributions reduction following Theorem 6]{JiaoHanWeissman2018}. The full-capacity coordinate, when included in the interval, is the unrestricted scalar oracle-value target studied by \citet{CausalSmithUnrestricted2026}.

The analysis connects constrained-allocation value processes \citep[Section 4 and Theorem 4.1.1]{FengHongNekipelov2026} and uniform inference for optimized value functions \citep{FirpoGalvaoParker2023} with large-alphabet functional estimation by polynomial approximation and moment matching \citep{JiaoVenkatHanWeissman2015,CaiLow2011}. Its risk criterion concerns the optimized population value across budgets; policy-learning work evaluates the welfare of an estimated assignment rule \citep{KitagawaTetenov2018,AtheyWager2021,Pellatt2023}. The formal statements and internal proofs are machine-checked in Lean 4, with verification scope recorded in the appendix.

The remainder first develops the connections to related work, then defines the causal class, budget value, and risk criterion. It next states the matched rate and summarizes the estimator, followed by discussion and limitations. The appendices give the supporting process bounds, lower-bound reduction, and verification scope.
